\section{Introduction}\label{sec:introduction}

\subsection{A first example}

Take the natural numbers with $+$, $\cdot$ and $/$, where division is undefined
exactly when the denominator is zero.  We want to total\-ize this partial
algebra without deciding what $0/0$ is.

The standard move is to keep the failed application as syntax.  Applications
that the base algebra can evaluate reduce to base values; the others remain
suspended terms.  So $6/2$ becomes the embedded natural $3$, while
\[
  0/0 \;=\; \susp(\mathsf{div},\old(0),\old(0))
\]
stays a term.  Nothing here is new: this is the free compatible completion, and
we use it only as a substrate.

Now ask a different question.  How much new structure must an observer add in
order to \emph{tell such terms apart}?

Fix a total algebra $T$ together with an injection of the base $\mathbb N$ into
it that fixes every natural pointwise, and require that $T$ add only
\emph{finitely many} states outside that copy of $\mathbb N$.  The base is
infinite, so $T$ is infinite; the budget counts only genuinely new semantic
states.  Call such a $T$ an \emph{observer}.

Two things happen, and they pull in opposite directions.

\begin{itemize}
  \item Given any finite list of terms, one observer separates all of them at
        once.  It keeps $\mathbb N$ pointwise, gives a distinct tag to each
        selected suspended subterm, and sends everything else to a single
        overflow state $\overflow$.
  \item Yet some sequences cannot be followed to a limit.  Start from $0/0$ and
        iterate the context $x\mapsto x+0$:
        \[
          0/0,\qquad 0/0+0,\qquad (0/0+0)+0,\qquad\ldots
        \]
        Because $a+0=a$ holds for every natural, this context fixes the entire
        base pointwise, so no observer can use base elements as memory for it.
        The terms grow syntactically forever, while each observer has only
        finitely many tags.  Given any candidate limit $x$, an observer can
        retain $x$ and a long prefix of the orbit exactly --- and then the next
        iterate must fall into $\overflow$, which is absorbing.  So $x$ is
        separated from every sufficiently late term of the sequence.
\end{itemize}

The sequence is therefore Cauchy for the separation these observers induce, and
it converges to nothing in the algebra.  The completion is \emph{proper}: it
has points that are not terms.  That is the phenomenon this paper is about.

\subsection{The general setting}

Let $D$ be a one-sorted binary partial algebra with carrier $A$.  We consider
compatible total extensions $T$ whose embedding of $A$ is injective, which fix
every element of $A$, and for which $|T\setminus A|$ is finite.  We write
$\overflow$ for the overflow state and $\stageeq{n}$ for indistinguishability
under all observers of budget $n$.

\begin{remark}[On the word \emph{observational}]
We use \emph{observational} in the metric sense: observers induce a separated
filtration and hence a completion.  This is not the observational
(behavioural) satisfaction of algebraic specification, where an observational
equivalence between algebras is induced by a distinguished set of observable
sorts \cite{BidoitHennicker2006,SannellaTarlecki2012}.  No quotient by an
equivalence is taken here; every base element is retained individually.
\end{remark}

The two mechanisms of the example are the two theorems of the paper.

\paragraph{Finite-pattern realization (\cref{thm:finite-pattern-realization}).}
For every finite family of generated Answers there is a single observer
injective on the whole family.  Pairwise separation is the two-point case, not
the primitive one.

\paragraph{Compression and escape (\cref{thm:compression-escape-properness}).}
Two independent facts combine:
\[
  \boxed{
  \text{finite trajectory compression}
  \;+\;
  \text{finite-pattern escape}
  \;\Longrightarrow\;
  \text{proper observational completion}.}
\]
Compression says that if the orbit relevant to a chosen unary family has a
uniformly finite code at each stage, then its factorial sample is Cauchy; this
is pigeonhole periodicity and factorial divisibility.  Escape says that no
generated Answer is a limit of a syntactically growing orbit, by the
candidate-tailored construction of the example.  Compression makes the sample
indistinguishable at every fixed stage; escape stops it converging.  Neither
alone gives properness.

\subsection{Why the observer class is the point}

The reader may expect the finite-family theorem to follow from the pairwise
case by taking products of separators, as in ordinary residual arguments.  It
does not, and the obstruction is visible already at the level of carriers.

Each one-point extension $\mathbb N\hookrightarrow\mathbb N\sqcup\{\overflow\}$
has a one-element complement.  But the diagonal copy of $\mathbb N$ inside the
product of two such extensions has infinitely many points $(a,\overflow)$
outside it.  The class ``fixes the base pointwise, finite complement'' is thus
\emph{not closed under binary products} over an infinite base
(\cref{prop:relative-product-obstruction}).  The finite-family theorem is
therefore proved by direct construction inside the class.

The same restriction is what makes escape non-trivial.  When the base is not
finitely codable, the candidate-tailored observer cannot itself be encoded into
any finite carrier, even though its external complement is finite.  It still
separates the candidate from the whole late tail.

\subsection{Quantitative consequence}

At stage $n$ the old-fixing orbit sees only $n+1$ external states --- the $n$
tags plus $\overflow$ --- unless it enters the pointwise-fixed base, after
which it stabilizes.  Hence for a distinct pair of iterates
\[
  c_{(n+1)!}\stageeq{n}c_{(n+2)!},
\]
and the least separation ranks are unbounded.  The factorial indices only
synchronize finite periods; the content is that finite compression and
finite-pattern escape coexist in the same observer class.

\subsection{Notation}

\begin{center}
\small
\begin{tabular}{@{}ll@{}}
\toprule
$D$, $A$ & partial algebra and its carrier; $A$ may be infinite\\
$\old(a)$ & the embedded copy of a base element $a\in A$\\
$\susp(f,s,t)$ & a suspended application: $f$ was undefined on $s,t$\\
$\RAlg{D}$ & the generated Answer algebra (free compatible completion)\\
observer & compatible total extension fixing $A$ pointwise, finite complement\\
$\overflow$ & the single overflow state of an observer\\
$\stageeq{n}$ & indistinguishable by every observer of budget $n$\\
$\Comp{\RAlg{D}}$ & the completion of the induced filtration\\
proper & the completion has points outside the image of $\RAlg{D}$\\
\bottomrule
\end{tabular}
\end{center}

\subsection{Organization and scope}

\Cref{sec:free-completion} records the substrate.
\Cref{sec:finite-separation} develops the observer class and proves
finite-pattern realization.
\Cref{sec:completion} builds the filtration and proves compression, escape and
their combination; the finite-base and old-fixing criteria follow as instances.
\Cref{sec:equations,sec:arithmetic} give exact equational conservativity and
the arithmetic examples.
\Cref{sec:diaconescu-bridge} locates the construction inside the standard
many-sorted partial-to-total encoding.
\Cref{sec:related-work,sec:limitations} position the work and list what is
open.

\paragraph{Scope.}
Free compatible completion, finite-state periodicity, factorial
synchronization, generic Cauchy completion, and abstract finite-set
discrimination are background.  What is offered here is the relative observer
class over a pointwise-fixed and possibly infinite base, the finite-pattern
realization theorem inside a class that is not product-closed, and the way
escape interacts with trajectory compression to force properness over infinite
bases.  Proofs are given in full in ordinary mathematical language; the
accompanying machine-checked development is described in
\cref{app:formalization} and is not relied on anywhere in the argument.
